\section{Introduction}\label{sec:intro}

\citet{krusellsmith1998} solve a heterogeneous-agent economy with aggregate risk by supposing that
households forecast prices from the aggregate shock and the mean of the wealth distribution, and
find that this forecast is almost exact. Behind the method lie questions the method does not
answer. Does an equilibrium exist? Is it recursive in some state, and in which? Is the economy
stationary? And when is the aggregate state $(z,K)$ --- shock and capital --- actually enough?

This note records what a Lean formalisation of the model proves about these questions. The
results are:

\begin{enumerate}[leftmargin=*,itemsep=2pt]
\item \textbf{Existence} (Section~\ref{sec:existence}). From every aggregate shock and initial
  distribution there is a sequential equilibrium, by Schauder--Tychonoff on paths of capital over
  the tree of shock histories; with logarithmic and CRRA utility; with Cobb--Douglas prices, the
  capital clamp provably slack under an explicit condition.
\item \textbf{Recursive equilibrium} (Section~\ref{sec:recursive}). With the value function, or
  the infinite history of aggregate shocks, added to the state, a recursive equilibrium exists in
  every economy; a stationary recursive equilibrium with an ergodic distribution exists in the
  sense of \citet{dgmm1994}, conditionally spotless in the sense of \citet{cao2020}. Uniqueness of
  sequential equilibrium implies a recursive equilibrium on the minimal state $(z,\mu)$.
\item \textbf{Uniqueness for small capital shares} (Section~\ref{sec:uniqueness}). For CRRA utility
  and Cobb--Douglas prices there is a capital share below which sequential equilibrium is unique,
  checked from primitives by a contraction in marginal utility. The law of motion then contracts in
  the distribution, which bounds the error of matching on capital \citep{lee2026}.
\item \textbf{Stationarity from the sequential equilibrium} (Section~\ref{sec:stationarity}). The
  equilibrium is a function of the infinite history of shocks, so it inherits stationarity from
  the shocks: the aggregate state of the unique small-share equilibrium has a unique, ergodic
  invariant distribution, with no recursive state space and no randomisation.
\item \textbf{The aggregate state $(z,K)$} (Section~\ref{sec:aggregate}). A recursive equilibrium on
  $(z,K)$ representing the equilibrium exists if and only if histories with the same $(z,K)$ have
  the same future capital and labour; then it is continuous and its stationary law is unique.
  Without idiosyncratic heterogeneity this always holds. The error of
  forecasting next capital from $(z,K)$ is at most the nonlinearity of saving plus the price
  feedback, times the distance between cross-sections.
\item \textbf{An explicit heterogeneous economy} (Section~\ref{sec:corner}). Log utility,
  Cobb--Douglas prices, employed households saving to a cap and unemployed households spending
  everything: sequential equilibrium is unique from every initial state, $(z,K)$ is sufficient, and
  the recursive equilibrium on $(z,K)$ has a unique, ergodic stationary distribution. Aggregate
  shocks have Krusell and Smith's persistence.
\item \textbf{Interior saving} (Section~\ref{sec:interior}). With prices depending on the shock
  only (a linear technology), i.i.d.\ shocks and cash ordered by the shock, $(z,K)$ is sufficient in
  every sequence-state equilibrium; in an explicit economy half of the households save strictly
  inside their constraints after every good shock.
\end{enumerate}

\noindent
Section~\ref{sec:open} lists what is open. Appendix~\ref{app:index} indexes the formal statements.

The prior literature is close in places. \citet{miao2006} proves existence of sequential equilibrium
by Schauder--Tychonoff and a recursive equilibrium with payoffs in the state; \citet{cao2020} proves
existence for the Krusell--Smith specification, a generalised recursive equilibrium, the value
function as extended state and, under uniqueness, the minimal-state recursive equilibrium, and
leaves the minimal-state question without uniqueness open; \citet{kublerschmedders2002} give
finite-agent examples where the wealth distribution is not a sufficient state. Both
\citet{miao2006} and \citet{kublerschmedders2002} treat uniqueness as a condition that cannot in
general be verified from primitives; Section~\ref{sec:uniqueness} verifies it from primitives, though
only for small capital shares. \citet{prohl2025} claims existence and uniqueness of a minimal-state recursive equilibrium by
monotone-operator methods; we found gaps in the proofs as written (the interiority argument in the
existence step, and convexity of the admissible set) and do not rely on it.
